% ratchet technical report (S. Kundu). Numbers: research/mismatch/results/kaggle-2026-10-15.txt and README.md in github.com/Shakhtar-Sankur/ratchet
\documentclass[11pt]{article}
\usepackage[final]{acl}
\usepackage{times}
\usepackage{latexsym}
\usepackage[T1]{fontenc}
\usepackage[utf8]{inputenc}
\usepackage{microtype}
\usepackage{courier}
\usepackage{booktabs}
\usepackage{amsmath}
\usepackage{xcolor}

\newcommand{\x}{$\times$}
\newcommand{\pending}[1]{\textcolor{red}{[PENDING: #1]}}

\title{The Sampler Was Not the Policy:\\Finding and Removing a Silent Training--Inference Mismatch\\in a From-Scratch GRPO Stack}

\author{Sankur Kundu \\
  Independent researcher, Bhubaneswar, India \\
  \texttt{sankur.kundu.tw@gmail.com}}

\begin{document}
\maketitle

\begin{abstract}
In reinforcement-learning post-training of language models, one engine samples answers and
another computes the gradient, and the policy-gradient estimate is only correct if both hold the
same policy. We build the whole loop from scratch, a C++/CUDA inference engine that returns the
log-probability of every sampled token and a PyTorch trainer in the engine's own weight layout,
so that every source of disagreement between them can be observed. At the released weights of
Qwen2.5-0.5B-Instruct the two agree to within 0.041 nats on 16,118 sampled tokens. After one or
two GRPO updates the gap reaches 5.2 nats and 0.65\% of tokens fall outside a 1$\pm$0.2
importance ratio: the engine stores weights in fp16, half an fp16 step around a typical weight
(7.6$\times$10$^{-6}$ at 0.02) exceeds an AdamW update at learning rate 10$^{-6}$, so the engine
rounds most of each update away and samples from a policy the trainer never had. Computing the
trainer's forward pass with the weights rounded exactly as the engine stores them, while the
optimizer keeps float32 master weights, brings the gap back to 0.043 with no token outside the
ratio bound. Taking the gap apart over 30 updates, rounding at weight sync accounts for nearly all
of it: 9--22\x\ the effect of fp16 arithmetic and 15--39\x\ that of the engine's own kernels. An
fp16 engine sees 5--13\% of weights change in an update; a bf16 engine 0.7--1.7\%, and its gap
grows to 11 nats with 15\% of tokens outside the ratio bound. With the fix, 100 GRPO steps on two
T4 GPUs raise GSM8K test accuracy from 48.1\% to 52.5\%. We also report the systems work that makes such a loop practical on small hardware:
in-place weight sync 369\x\ faster than a reload, partial rollouts with one-step-ahead training
that make steps 1.9\x\ faster, and a sort-free sampler that makes rollouts 27\x\ faster. Code and
raw logs are released.\footnote{\url{https://github.com/Shakhtar-Sankur/ratchet}}
\end{abstract}

\section{Introduction}

Policy-gradient post-training of language models, as in PPO \citep{schulman2017ppo} and GRPO
\citep{shao2024deepseekmath}, alternates between sampling answers and learning from them. At any
scale the two are done by different software: a serving engine optimised for generation
\citep{kwon2023vllm} and a training framework optimised for gradients
\citep{sheng2024hybridflow,hu2024openrlhf}. Each computes the probability of the same token in
its own way, and the objective's importance ratio is assumed to be one when their weights are
equal. When it is not, the gradient is biased and nothing announces it.

Recent work measures this gap at the level of logits and corrects for it, for example by
switching the training precision from bf16 to fp16 \citep{qi2025fp16}. We take a complementary,
bottom-up approach: we control both sides completely, so each source of disagreement can be
switched on alone and measured on the same sampled tokens. This report makes three contributions.

\begin{enumerate}
\item \textbf{A from-scratch, observable RL loop} (\S\ref{sec:system}): an inference engine and a
trainer that share a weight layout, agree to within 10$^{-6}$ on the CPU, and log the gap on every
step of every run.
\item \textbf{A silent failure, its cause and its fix} (\S\ref{sec:mismatch},
\S\ref{sec:decompose}): weight sync into an fp16 engine rounds away most of each update, and this
rounding, not arithmetic precision or kernels, is where the gap comes from; computing the trainer
with the engine's rounded weights, as mixed-precision training does \citep{micikevicius2017mixed},
removes it.
\item \textbf{Systems measurements on two T4 GPUs} (\S\ref{sec:systems}, \S\ref{sec:results}):
weight sync, the long tail of answer lengths, sampling cost, and end-to-end accuracy, including
the results that are within noise.
\end{enumerate}

\section{System}
\label{sec:system}

\paragraph{Engine.} Rollouts are generated by relay, a C++/CUDA inference engine with paged KV
cache and continuous batching, through a C interface. Every sampled token comes with log softmax
of the logits divided by the temperature, over the whole vocabulary: the behaviour policy's
log-probability. A rollout stopped after any number of tokens and resumed gives the same tokens
and log-probabilities as one that ran through, which partial rollouts rely on.

\paragraph{Trainer.} The policy is a Llama/Qwen2 decoder in PyTorch whose parameters have the
engine's names and fused layout, so a weight sync is a copy with no renaming or reshaping. On four
model shapes (multi-head, grouped-query with tied bf16 weights, Qwen2 biases, Llama~3 RoPE
scaling) its logits match the \texttt{transformers} reference exactly or within
2.1$\times$10$^{-7}$, and its log-probabilities match the engine's rollouts within
9.5$\times$10$^{-7}$ on the CPU. A whole run is reproducible bit for bit.

\paragraph{GRPO.} Each prompt is sampled eight times; rewards are 1 when the final number of the
answer equals the reference \citep{cobbe2021gsm8k}; advantages are normalised within each group;
the loss is PPO's clipped objective averaged over response tokens, optimised with AdamW. Groups
whose answers all received the same reward are skipped.

\section{The mismatch}
\label{sec:mismatch}

\paragraph{At the released weights, the two agree.} On a T4 with Qwen2.5-0.5B-Instruct
\citep{qwen2025qwen25}, the engine runs in fp16 and the trainer in float32. Over 16,118 sampled
tokens the absolute difference between their log-probabilities is at most 0.041 (mean
6.8$\times$10$^{-4}$, 99th percentile 0.0085), and no token's importance ratio lies outside
1$\pm$0.2. The released weights are bf16 values that the engine's fp16 storage represents, so at
this point both sides hold the same weights.

\paragraph{After an update, they do not.} After one or two GRPO steps, the gap on fresh rollouts
reached 4.5 to 6.4 nats on some tokens. The cause is in the weight sync, not the arithmetic. An
AdamW step at learning rate 10$^{-6}$ moves a weight by roughly 10$^{-6}$; for a weight of 0.02,
adjacent fp16 values are 1.5$\times$10$^{-5}$ apart, so any change smaller than half of that,
7.6$\times$10$^{-6}$, rounds back to the old value when the engine converts the trainer's float32
weights. The engine therefore keeps the update only on small weights and samples from a policy
that the trainer never had: the trainer computes ratios and gradients for one policy, while the
samples come from another.

\paragraph{The fix.} The trainer computes its forward pass with each matrix rounded to fp16 exactly
as the engine stores it (round to nearest even), and the optimizer updates float32 master weights
through a straight-through gradient, which is the standard mixed-precision recipe
\citep{micikevicius2017mixed} applied across the trainer--engine boundary. Small updates now
accumulate in the master weights until they cross an fp16 step, and both sides always compute with
the same values.

\begin{table}[t]
\centering
\small
\begin{tabular}{lrrrr}
\toprule
trainer weights & max & mean & p99 & ratio off \\
\midrule
float32 & 5.21 & 0.0117 & 0.155 & 0.65\% \\
fp16-rounded & \textbf{0.043} & \textbf{0.00067} & \textbf{0.0078} & \textbf{0} \\
\bottomrule
\end{tabular}
\caption{Absolute log-probability gap between engine and trainer on the same 64 fresh rollouts
(19,313 tokens) after two GRPO steps. Ratio off: share of tokens whose importance ratio
leaves 1$\pm$0.2.}
\label{tab:gap}
\end{table}

\paragraph{Result.} Table~\ref{tab:gap}. With fp16-rounded weights the gap after training is what it
was at the released weights. In every later run, colocated and split, the per-step maximum stayed
between 0.02 and 0.06.

\section{Systems}
\label{sec:systems}

\paragraph{Weight sync.} The simple sync copies every parameter to the host, converts it to fp16 on
the CPU and rebuilds the engine. Instead, the trainer hands the engine each parameter's device
pointer; the engine converts float32 to fp16 on its own GPU, in place, producing the same bits as
the host path, and drops its prefix cache, whose entries were computed with the old weights.
Updates are refused while a rollout is running, and a later full reload is refused rather than
silently reverting the policy. On one T4 with 494M parameters (1.98\,GB in float32), the sync
takes 0.017\,s instead of 6.26\,s (369\x); across two GPUs, with a peer copy, 0.24\,s. Rollouts
after either sync are identical.

\paragraph{The long tail.} A step waits for its longest answer. Partial rollouts
\citep{kimi2025k15} keep more groups generating than a step trains on, end the step when enough
have finished, and resume the rest under the new weights from the exact token where they stopped.
One-step-ahead training generates batch $k{+}1$ while batch $k$ trains, as asynchronous systems do
\citep{fu2025areal}. Both train on tokens sampled by an older policy, so the importance ratio
always uses the engine's sampling-time log-probabilities. On two T4s, one GPU generating while the
other trains, a synchronous step takes 49.1\,s and a step with both techniques 25.6\,s (1.9\x),
about the training time alone.

\paragraph{Sampling.} The first GPU runs generated 34 tokens per second: the engine sorted all
151,936 vocabulary entries on the CPU for every token, which took longer than the forward pass.
Sampling at temperature 1 with no top-$k$ or top-$p$, as RL does, needs no sort; walking the
inverse cumulative distribution in token order, with the batch's rows sampled in parallel, takes
1.3\,ms per token instead of 29.8\,ms, and rollouts run at 931 tokens per second (27\x).

\section{Results}
\label{sec:results}

\begin{table}[t]
\centering
\small
\begin{tabular}{lrrr}
\toprule
mode & before & after & time \\
\midrule
colocated & 48.1\% & \textbf{52.5\%} & 50.6 min \\
split, ahead + partial & 47.8\% & 50.2\% & 44.7 min \\
\bottomrule
\end{tabular}
\caption{Greedy accuracy on all 1,319 GSM8K test problems before and after 100 GRPO steps (8
problems $\times$ 8 answers per step, learning rate 10$^{-6}$), Qwen2.5-0.5B-Instruct, two T4s.
Colocated: both GPUs generate, then train with data parallelism. Split: one generates while the
other trains. One run per mode.}
\label{tab:gsm8k}
\end{table}

Table~\ref{tab:gsm8k}. Training reward rose from 0.44 to 0.58 (colocated, first against last ten
steps) and from 0.38 to 0.56 (split); answers cut off at 512 tokens fell from 6\% to 2\%. The same
weights scored 48.1\% and 47.8\% under two batchings of the test set, because fp16 results depend
slightly on which sequences share a batch, so differences of a few tenths of a point are noise.
With one run per mode, the difference between the two gains is not established.

\section{Decomposing the gap by source}
\label{sec:decompose}

Because both sides are ours, each source of disagreement can be switched on alone and measured on
the same sampled tokens. \emph{Weight rounding} compares the trainer with its float32 master
weights against the same trainer with those weights rounded to fp16, as the engine stores them, or
to bf16. \emph{Arithmetic precision} compares fp16 autocast with float32 at identical weights.
\emph{Engine kernels} is what remains between the engine and the closest trainer variant. Each is
measured on 64 fresh rollouts (8 problems $\times$ 8 answers, temperature 1) at the released
weights and after 1, 3, 10 and 30 GRPO updates of a run with the fix, on two T4s. We also count,
after every update, the share of weight-matrix entries whose rounded value changes: the part of the
update the engine can see.

\begin{table*}[t]
\centering
\small
\begin{tabular}{rrrrrrrrrrr}
\toprule
& \multicolumn{3}{c}{weight rounding, fp16} & \multicolumn{3}{c}{weight rounding, bf16} & \multicolumn{2}{c}{arithmetic} & \multicolumn{2}{c}{engine kernels} \\
\cmidrule(lr){2-4}\cmidrule(lr){5-7}\cmidrule(lr){8-9}\cmidrule(lr){10-11}
updates & mean & max & off & mean & max & off & mean & max & mean & max \\
\midrule
0 & 0.000 & 0.0001 & 0 & 0 & 0 & 0 & 0.0013 & 0.081 & 0.0007 & 0.040 \\
1 & 0.016 & 1.93 & 1.2\% & 0.022 & 2.16 & 2.2\% & 0.0012 & 0.073 & 0.0007 & 0.033 \\
3 & 0.025 & 2.30 & 2.8\% & 0.041 & 3.65 & 5.5\% & 0.0012 & 0.136 & 0.0007 & 0.056 \\
10 & 0.030 & 3.47 & 3.4\% & 0.104 & 9.70 & 13.3\% & 0.0014 & 0.064 & 0.0008 & 0.031 \\
30 & 0.012 & 0.78 & 1.0\% & 0.196 & 10.98 & 15.0\% & 0.0013 & 0.069 & 0.0008 & 0.066 \\
\bottomrule
\end{tabular}
\caption{The engine--trainer gap taken apart: absolute difference in per-token log-probability
(nats) caused by each source alone, on 64 fresh rollouts per checkpoint of one run. Off: share of
tokens whose importance ratio leaves 1$\pm$0.2 (none for arithmetic or kernels at any checkpoint).}
\label{tab:decompose}
\end{table*}

\paragraph{Rounding at sync is the gap.} At the released weights every source is small, and
rounding is zero for bf16 because the released weights are bf16 values. After a single update,
rounding the same weights to fp16 changes per-token log-probabilities by 0.016 on average and by up
to 1.9 nats, with 1.2\% of tokens outside the ratio bound (Table~\ref{tab:decompose}). Arithmetic
precision and the engine's kernels stay where they were at every checkpoint, at most 0.14 and 0.07
nats with no token outside the bound; on average, rounding is 9--22\x\ larger than arithmetic and
15--39\x\ larger than the kernels. The engine is always closest to the trainer variant that uses
its own rounded weights, which is the fix of \S\ref{sec:mismatch}: with it, the engine and trainer
agree within 0.031--0.066 nats at every checkpoint. A second run without the fix reached 1.8 to
4.6 nats between engine and trainer by updates 1--10, as in \S\ref{sec:mismatch}.

\paragraph{How much of an update an engine sees.} Once rounded to fp16, only 13.4\% of
weight-matrix entries change value in the first update, falling to 5.0\% by update 30 as the
optimizer's steps shrink; with the fix, the rest of each update accumulates in the master weights
until it crosses an fp16 step. In bf16, which has three fewer mantissa bits, 1.7\% change in the
first update and 0.7\% by update 30, about eight times fewer. These shares count every entry,
including any the update did not move, such as embedding rows of tokens absent from the batch. From
the first update on, 97--99.8\% of master weights lie off either grid.

\paragraph{bf16 weights would be worse, and grow worse.} Had the engine stored bf16 weights, the gap from
rounding alone would have been larger at every checkpoint and grown with training, to a mean of
0.20 and a maximum of 11 nats, with 15\% of tokens outside the ratio bound by update 30, while the
fp16 gap stayed between 0.012 and 0.030. This agrees with the finding that fp16 reduces the
training--inference mismatch relative to bf16 \citep{qi2025fp16}; here the cause can be seen
directly, in the weights the engine holds. Why the bf16 gap grows in this run is not established by
one run.

\paragraph{Where the gap falls.} Without the fix, the gap concentrates on the tokens the engine
sampled with low probability: a mean of 0.10--0.17 nats for tokens it gave less than 0.1, against
0.001--0.002 for tokens above 0.9. These are the exploratory tokens whose importance ratios the
clipped objective acts on.

\section{Discussion}

The failure in \S\ref{sec:mismatch} is easy to miss: training still runs, rewards still move, and
the only symptom is a gap that most stacks do not measure. Section~\ref{sec:decompose} shows where
to look: not at the arithmetic or the kernels, which stay within 0.14 nats, but at the
weights each side actually holds after a sync. It is also easy to cause, because
storing inference weights in a lower precision than the optimizer's is common in systems that
separate the two. Two habits catch it: logging the engine--trainer gap on fresh
samples at every step, and computing the trainer's forward pass with the weights the engine will
actually hold. The second costs one rounding per matrix per step.

\section*{Limitations}

All results use one 0.5B model, one task and two T4 GPUs, with one run per training mode. The
decomposition in \S\ref{sec:decompose} uses 64 rollouts per checkpoint from one run; its bf16
column measures rounding to bf16 in our trainer, not a bf16 engine. Larger models and other
learning rates change how much of an update survives rounding.

\section*{Use of AI assistants}

AI coding and writing assistants were used to help write code and prose. The author designed the
study, checked every reported number against the released logs, and takes responsibility for the
content.

\bibliography{refs}

\end{document}
